% ── Chapter 4: Design & Architecture ─────────────────────────────────────────
\chapter{Design and Architecture}

\section*{Introduction}
\addcontentsline{toc}{section}{Introduction}

Chapter~3 closed on a clear diagnosis: the information needed to steer the business exists, but it
lies scattered across isolated operational applications, inconsistent and hard to exploit, leaving
management to decide on intuition rather than on a reliable, consolidated reading of the activity.
This chapter is the pivot from that problem to its solution. It presents LOGIQ,
the Business Intelligence decision-support system proposed to close the analytical gap, and sets out
\emph{what} the system must do and \emph{how} it is logically structured to do it.

The chapter is deliberately \emph{conceptual}. It stays at the level of functions and logical
structure --- the actors and capabilities the system supports, the modules it decomposes into, the
layered architecture that organises it, the way the analytical data is modelled, the integration
process that feeds it, and the application through which decision-makers consult it. Concerns that
belong to construction --- the technology stack, the frameworks, the physical deployment --- are
deliberately deferred to Chapter~5, which mirrors this design as its technical realisation. The aim
here is a faithful blueprint, independent of any particular implementation.

The presentation follows a single thread, from the abstract to the concrete. It begins with the
\emph{functional specification} (what the system must do, traced back to the objectives of
Chapter~3), then derives the \emph{functional architecture} (how those functions are decomposed and
related), and the \emph{technical specifications} that frame its realisation. From there it introduces
LOGIQ and its \emph{global architecture}, before zooming in on
the three pillars that give the system its value: the \emph{data warehouse} at its core, the
\emph{ETL pipeline} that populates it, and the \emph{application platform} through which it is used.
The chapter closes on the cross-cutting concerns of \emph{security and conformity}, and on a
transition towards implementation.

% ═════════════════════════════════════════════════════════════════════════════
\section{Functional Specification}
\label{sec:ch4-functional}

Before describing how LOGIQ is structured, this section sets out \emph{what} it must do. In keeping
with the spirit of Chapter~3, the specification is expressed not as a formal requirements list but as
a \emph{functional design}: the actors that interact with the system, the capabilities it offers
them, and the use cases through which those capabilities are exercised. Every capability is traced
back to one of the specific objectives derived in Chapter~3 (\S\ref{sec:ch3-objectives}), so that the
design answers the diagnosed needs rather than adding features for their own sake.

% -----------------------------------------------------------------------------
\subsection{Actors and Roles}

LOGIQ is used not by the public but by Yalidine's own decision-makers, each consulting the system
from the angle of their responsibility. Users are not self-registered: they are imported from the
company's HR system and activated by an administrator, who assigns each one a \textbf{role}. The role
then governs which views the user may consult --- the foundation of the role-appropriate access
objective. Table~\ref{tab:ch4-actors} lists the \emph{default} roles the system ships with and the
views each one consults.

% [VISUAL-IMPLEMENT] default roles of the system and the views each one consults
\begin{table}[H]
\centering
\renewcommand{\arraystretch}{1.35}
\rowcolors{2}{white}{tblrowalt}
\begin{tabularx}{\linewidth}{|>{\raggedright\arraybackslash}p{4.3cm}|>{\raggedright\arraybackslash}X|}
\hline
\rowcolor{tblheader}
\textcolor{white}{\textbf{Default role}} & \textcolor{white}{\textbf{What they consult}} \\
\hline
Direction Générale \newline (general management) &
Consolidated overview of both axes \newline
Parcel Delivery view (Operations, Cost \& Profitability, Performance) \newline
Dedicated Transport view (Operations, Cost \& Profitability, Performance) \newline
Route Analysis view (future work) \\
Responsable Colis \& PCC \newline (parcel \& PCC manager) &
Consolidated overview \newline
Parcel Delivery view (Operations, Cost \& Profitability, Performance) \\
Responsable Tournées \newline (routes manager) &
Consolidated overview \newline
Route Analysis view (future work) \\
Responsable Transport \newline (transport manager) &
Consolidated overview \newline
Dedicated Transport view (Operations, Cost \& Profitability, Performance) \\
Administrator &
User activation \newline
Role assignment \newline
Data-pipeline run history \\
\hline
\end{tabularx}
\caption{The default roles of LOGIQ and the views each one consults.}
\label{tab:ch4-actors}
\end{table}

The roles in Table~\ref{tab:ch4-actors} are only the \emph{default} roles shipped with the system.
Because access is governed by configurable, view-level permissions, the administrator can create and
manage additional roles --- granting or restricting access to any view --- so that new
decision-makers can be onboarded as the organisation evolves, without changing the system.

% -----------------------------------------------------------------------------
\subsection{Functional Specifications}

The specifications expected of the system group into five families, each answering one or more of the
specific objectives of Chapter~3 (Table~\ref{tab:ch4-capabilities}). They run from bringing the
scattered data together, through consulting it along the two axes and three views, to being warned
proactively when an indicator drifts --- the whole governed by role-appropriate access and supported
by administration.

% [VISUAL-IMPLEMENT] functional specification families traced to the Ch.3 objectives
\begin{table}[H]
\centering
\renewcommand{\arraystretch}{1.35}
\rowcolors{2}{white}{tblrowalt}
\begin{tabularx}{\linewidth}{|>{\raggedright\arraybackslash}p{3.0cm}|>{\raggedright\arraybackslash}X|>{\raggedright\arraybackslash}p{3.3cm}|}
\hline
\rowcolor{tblheader}
\textcolor{white}{\textbf{Specification family}} & \textcolor{white}{\textbf{What the system must do}} &
\textcolor{white}{\textbf{Objective served (Ch.3)}} \\
\hline
Data consolidation &
Bring the scattered operational data into a single, coherent source \newline
Refresh it continuously &
Consolidate the scattered data (Obj.~1). \\
Multi-axis decision consultation &
Consult each business axis --- parcel delivery and dedicated transport --- along its three views (Operations, Cost \& Profitability, Performance) \newline
Show a consolidated overview of both axes &
Read each axis in depth (Obj.~2, 3). \\
Proactive monitoring &
Define threshold rules on sensitive indicators \newline
Alert automatically, through the preferred channel, when a metric drifts beyond its acceptable level &
Monitor proactively (Obj.~4). \\
Role-appropriate access &
Present each decision-maker only the sections suited to their role \newline
Let them save personal filters and notification preferences &
Deliver the right information to the right decision-maker (Obj.~5). \\
Administration &
Activate HR-imported users \newline
Assign and adjust roles \newline
Supervise the data-refresh pipeline and its run history &
Underpin reliable, governed operation; found decisions on reliable data (Obj.~1, 6). \\
\hline
\end{tabularx}
\caption{The five functional specification families, each traced to a specific objective of Chapter~3.}
\label{tab:ch4-capabilities}
\end{table}

% -----------------------------------------------------------------------------
\subsection{Use-Case Overview}

Figure~\ref{fig:ch4-usecase} draws actors and use cases together as a use-case overview. The
business user interacts mainly with the consultation and monitoring use cases, while the
administrator governs users, roles, and the data pipeline. Access to every use case is mediated by
the user's role, so two users signing in under different roles may legitimately see different
sections.

% [VISUAL-IMPLEMENT] global use-case diagram (actors x main use cases, role-mediated)
% Stick-figure actors + ellipse use cases inside a LOGIQ system boundary; palette from settings.tex.
\begin{figure}[H]
\centering
\begin{tikzpicture}[
  >={Stealth[length=2.4mm]},
  uc/.style={draw=figblue, fill=figbluetint, ellipse, align=center, text=black!75,
             minimum width=3.5cm, minimum height=1.0cm, inner sep=1pt, font=\scriptsize},
  rel/.style={draw=figslate, line width=0.6pt},
  alab/.style={font=\scriptsize\bfseries, align=center, text=black!70},
  % reusable stick-figure actor
  actor/.pic={
    \draw[line width=0.7pt, figslate] (0,0.55) circle (0.16);
    \draw[line width=0.7pt, figslate] (0,0.39) -- (0,-0.10);
    \draw[line width=0.7pt, figslate] (-0.22,0.20) -- (0.22,0.20);
    \draw[line width=0.7pt, figslate] (0,-0.10) -- (-0.18,-0.42);
    \draw[line width=0.7pt, figslate] (0,-0.10) -- (0.18,-0.42);
  }]

  % ---- system boundary -------------------------------------------------------
  \node[draw=figslate, line width=0.8pt, rounded corners=4pt, fill=black!2,
        minimum width=4.8cm, minimum height=9.6cm,
        label={[font=\bfseries, text=figblue]above:LOGIQ}] (sys) at (0,-0.4) {};

  % ---- use cases (centre column) ---------------------------------------------
  \node[uc] (uc1) at (0, 3.5)  {Consult consolidated\\overview};
  \node[uc] (uc2) at (0, 1.9)  {Analyse an axis\\{\itshape Operations · Cost \&}\\{\itshape Profitability · Performance}};
  \node[uc] (uc3) at (0, 0.2)  {Manage \& receive\\proactive alerts};
  \node[uc] (uc4) at (0,-1.4)  {Manage filters\\\& preferences};
  \node[uc] (uc5) at (0,-3.0)  {Administer\\users \& roles};
  \node[uc] (uc6) at (0,-4.5)  {Supervise the\\data pipeline};

  % ---- actor (left): the business user --------------------------------------
  \pic at (-5.4, 1.0)  {actor};
  \node[alab] at (-5.4, 0.18) {Business user\\{\itshape (utilisateur métier)}};

  % ---- actor (right) ---------------------------------------------------------
  \pic at (5.4,-3.0)  {actor};
  \node[alab] at (5.4,-3.70) {Administrator};

  % ---- associations: business user -> use cases ------------------------------
  \draw[rel] (-5.1, 1.0) -- (uc1.west);
  \draw[rel] (-5.1, 1.0) -- (uc2.west);
  \draw[rel] (-5.1, 1.0) -- (uc3.west);
  \draw[rel] (-5.1, 1.0) -- (uc4.west);

  % ---- associations: administrator -> use cases ------------------------------
  \draw[rel] (5.1,-3.0) -- (uc5.east);
  \draw[rel] (5.1,-3.0) -- (uc6.east);
  \draw[rel] (5.1,-3.0) -- (uc3.east);
\end{tikzpicture}
\caption{Global use-case overview of LOGIQ: the business user consults and monitors; the administrator
governs users, roles, and the pipeline. Every association is mediated by the user's role.}
\label{fig:ch4-usecase}
\end{figure}

% -----------------------------------------------------------------------------
\subsection{Quality Expectations}
\label{subsec:ch4-quality}

Beyond what the system does, Chapter~3 attached three quality conditions without which it would not be
adopted (Table~\ref{tab:ch3-needs}). They are restated here as functional-level expectations that
shape the design, rather than as a formal non-functional specification:

\begin{itemize}[noitemsep]
    \item \textbf{Reliability} : the figures presented must be trustworthy and consistent across
          views; a decision-support system is only as valuable as the confidence its users place in
          its numbers.
    \item \textbf{Responsiveness at national scale} : the system must stay fast over data covering
          Yalidine's whole national activity, so that opening a view or reading a long history returns
          in seconds, not minutes.
    \item \textbf{Simplicity for non-technical managers} : the interface must be readable and
          navigable by managers who are not data specialists, favouring clear indicators and visual
          reading over raw tables.
\end{itemize}

These three expectations recur as design drivers throughout the chapter: reliability motivates the
layered, quality-gated warehouse; responsiveness motivates the pre-computed aggregate layer and the
separation of analytical from operational data; and simplicity motivates the dashboard organisation
around two axes and three readable views.


% ═════════════════════════════════════════════════════════════════════════════
\section{Functional Architecture}
\label{sec:ch4-funcarch}

The functional specification said \emph{what} LOGIQ must do; the functional architecture says \emph{how
those functions are organised} into a coherent whole --- still at a purely logical level, before any
technology is chosen. Read left to right, the system is a single, one-directional \emph{value chain}
that turns raw activity into a decision: data leaves the operational systems, is \textbf{processed},
\textbf{consolidated} into one reliable base, and finally delivered to the decision-maker as
\textbf{presentation} and \textbf{alerts}. Figure~\ref{fig:ch4-funcarch} shows this chain and the
modules that make it up.

% [VISUAL-FETCH] functional-architecture block diagram (designed externally; tech-agnostic value chain)
\fig[\textwidth]{functional-architecture}{Functional architecture of LOGIQ: a one-directional value
chain from the heterogeneous operational data sources, through data processing and reliable
consolidation, to the two decision-ready outputs --- presentation and alerts --- consumed by the
decision-maker.}{ch4-funcarch}

The chain begins outside LOGIQ, in the \textbf{data sources and operational applications} --- the
several independent systems that run the business day to day and in which its data lies scattered.
\textbf{Data Processing} draws from them and brings the data into a consistent form, which
\textbf{Reliable Data Consolidation} then holds as a single, coherent, continuously refreshed base ---
the common ground every other module reads, and the point at which the figures become trustworthy.
From that base the chain forks into the two \emph{decision-ready outputs}: \textbf{Data Presentation},
which the decision-maker reads on demand as the two axes and their three views, and \textbf{Alerts \&
Notifications}, which watches the same indicators and pushes a warning the moment one drifts.
Presentation answers questions when they are asked; alerting answers them before they are asked.

Two further responsibilities do not sit at a single point in the chain but cut across all of it, and
are left out of Figure~\ref{fig:ch4-funcarch} for clarity: \textbf{Access \& Personalisation}, which
filters every interaction to the user's role and remembers their saved filters and notification
preferences, and \textbf{Administration}, which governs the users, the roles, and the processing
pipeline that keeps the base trustworthy. Table~\ref{tab:ch4-funcmodules} ties every module --- those
in the figure and these two cross-cutting concerns --- back to the specification family it realises.

% [VISUAL-IMPLEMENT] functional modules mapped to the specification families of S4.1
\begin{table}[H]
\centering
\renewcommand{\arraystretch}{1.35}
\rowcolors{2}{white}{tblrowalt}
\begin{tabularx}{\linewidth}{|>{\raggedright\arraybackslash}p{3.6cm}|>{\raggedright\arraybackslash}p{3.0cm}|>{\raggedright\arraybackslash}X|}
\hline
\rowcolor{tblheader}
\textcolor{white}{\textbf{Functional module}} & \textcolor{white}{\textbf{Realises}} &
\textcolor{white}{\textbf{Conceptual role}} \\
\hline
Data Processing & Data consolidation & Draw data from the heterogeneous operational sources and bring
it into a consistent form. \\
Reliable Data Consolidation & Data consolidation & Hold the processed data as one coherent,
continuously refreshed and trustworthy analytical base. \\
Data Presentation & Multi-axis decision consultation & Expose the two business axes along Operations /
Cost \& Profitability / Performance, plus the consolidated overview of both. \\
Alerts \& Notifications & Proactive monitoring & Watch sensitive indicators and push an alert through
each user's preferred channel when one drifts. \\
Access \& Personalisation \newline \emph{(cross-cutting)} & Role-appropriate access & Mediate every
interaction by role, and hold each user's saved filters and notification preferences. \\
Administration \newline \emph{(cross-cutting)} & Administration & Activate HR-imported users, manage
roles, and supervise the processing pipeline and its run history. \\
\hline
\end{tabularx}
\caption{The functional modules of LOGIQ, each realising one specification family of
\S\ref{sec:ch4-functional}; the last two cut across the chain and are not shown in
Figure~\ref{fig:ch4-funcarch}.}
\label{tab:ch4-funcmodules}
\end{table}

This decomposition is the bridge between the specification and the concrete solution: each module is a
\emph{responsibility}, not yet a component. The next section introduces LOGIQ and shows how these
responsibilities are projected onto a layered, service-oriented architecture.

% ═════════════════════════════════════════════════════════════════════════════
\section{Technical Specifications}
\label{sec:ch4-techspec}

The quality expectations of \S\ref{subsec:ch4-quality} stated, in the decision-maker's own terms, the
conditions under which LOGIQ would be trusted and adopted. This section turns them --- together with
the engineering constraints they imply at national scale --- into a structured \emph{technical
specification}: the non-functional properties the design must guarantee, and against which its
realisation in Chapter~5 can be judged. To keep the specification systematic, it is organised along
the product-quality characteristics of the ISO/IEC~25010 model~\cite{iso25010}; and, like the rest of
this chapter, each requirement is stated independently of any particular technology.
Table~\ref{tab:ch4-techspec} collects them.

% [VISUAL-IMPLEMENT] technical (non-functional) specification, organised by ISO/IEC 25010 attributes
\begin{table}[H]
\centering
\renewcommand{\arraystretch}{1.35}
\rowcolors{2}{white}{tblrowalt}
\begin{tabularx}{\linewidth}{|>{\raggedright\arraybackslash}p{3.0cm}|>{\raggedright\arraybackslash}X|>{\raggedright\arraybackslash}p{2.8cm}|}
\hline
\rowcolor{tblheader}
\textcolor{white}{\textbf{Quality attribute}} & \textcolor{white}{\textbf{Technical specification}} &
\textcolor{white}{\textbf{Driven by}} \\
\hline
Performance efficiency & Analytical views and KPI queries return within a few seconds --- not minutes
--- over the full national history; response time and resource use stay bounded as the data grows, by
reading from a pre-computed aggregate layer and isolating analytical from operational workloads. &
Responsiveness \\
Capacity \& scalability & Sized for production-scale volumes (a parcel history on the order of hundreds
of millions of status events over a multi-year horizon), refreshed incrementally rather than reloaded
in full, leaving headroom to grow with the activity. & National scale \\
Reliability & Figures are consistent across every view; the data refresh is idempotent and
recoverable, so a failed or repeated run never corrupts or duplicates data; the system degrades
gracefully when the analytical store is briefly unavailable. & Reliability \\
Security & Every view is gated by role-based access control; users are provisioned only by an
administrator; the two financial perimeters are kept strictly separate; sessions use stateless
authentication and all transport is encrypted. & Role-appropriate access \\
Usability & Readable and navigable by managers who are not data specialists --- visual-first, clear
indicators over raw tables --- and consistent across the two axes and their three views. &
Simplicity \\
Maintainability & Built as the loosely coupled functional modules of \S\ref{sec:ch4-funcarch} on open,
non-proprietary foundations, with a clean separation of operational and analytical concerns, so each
part can evolve independently. & Sustainable evolution \\
Compatibility \& portability & Source integration depends only on each system's API contract --- so the
real production systems connect unchanged --- and the design accommodates the deferred Route Analysis
axis without restructuring; the whole is reproducibly deployable on modest infrastructure. &
Integration \& future work \\
\hline
\end{tabularx}
\caption{Technical specification of LOGIQ, organised along the product-quality characteristics of
ISO/IEC~25010 and traced to the quality expectations of \S\ref{subsec:ch4-quality}.}
\label{tab:ch4-techspec}
\end{table}

Three of these characteristics dominate the design and recur as drivers in the sections that follow:
\textbf{performance and capacity} justify the pre-computed aggregate layer and the analytical/operational
split of the warehouse; \textbf{reliability} justifies its layered, quality-gated construction and an
idempotent refresh; and \textbf{security} justifies role-based access and the strict separation of the
two financial perimeters. The remaining characteristics --- usability, maintainability, and
portability --- shape, respectively, the dashboard, the modular decomposition just presented, and the
deployment, and are revisited where each is realised.

% ═════════════════════════════════════════════════════════════════════════════
\section{Solution and Global Architecture}
% TODO

% ═════════════════════════════════════════════════════════════════════════════
\section{Data Warehouse Design}
% TODO

% ═════════════════════════════════════════════════════════════════════════════
\section{ETL Pipeline Design}
% TODO

% ═════════════════════════════════════════════════════════════════════════════
\section{Application Platform Design}
% TODO

% ═════════════════════════════════════════════════════════════════════════════
\section{Security and Conformity}
% TODO

% ═════════════════════════════════════════════════════════════════════════════
\section*{Conclusion}
\addcontentsline{toc}{section}{Conclusion}
% TODO
